\documentclass[11pt]{article}
\usepackage{amsmath,amssymb,amsthm,booktabs}
\usepackage[margin=2.4cm]{geometry}
\usepackage[hyphens]{url}\sloppy
\newtheorem{theorem}{Theorem}\newtheorem{lemma}[theorem]{Lemma}\newtheorem{proposition}[theorem]{Proposition}
\newtheorem{corollary}[theorem]{Corollary}
\theoremstyle{definition}\newtheorem{definition}[theorem]{Definition}\newtheorem{remark}[theorem]{Remark}
\newcommand{\Z}{\mathbb Z}
\title{Q597 case a: (Ea) and (Stale$_a$) are proved, so (Ca) and the case-a half of BS$(L)$ hold\\
\normalsize PROVE 2026-11-23 c2. In case~a ($L$'s re-run takes the other class-1 a-step) both hypotheses of Proposition~casea (11-22 c2) are now theorems. \textbf{(Ea)} has a short proof. The engine stops a journey toward $T$ only on nest $T$'s two-dimensional cell grid, which is the common level set of two linear functionals, one on each lattice. Both functionals are monotone along every $T$-journey, and each moves strictly at exactly one of the two class-1 a-steps. So the other journey's first middle step certifies that $R_p$ is off the grid in the coordinate that the first journey's step leaves unchanged. \textbf{(Stale$_a$)} is proved in two parts. First, the board difference at $p+2$ is the remaining wake together with the non-rhombus tiles that $A$ made at $p,p+1$: the re-run's two hexagons read, outside $A$'s, only $J_o$-created tiles (finite check). Second, $A$ reads no wake tile after $p+1$. That is proved by hand from Lemma~phi (future and contact steps of $J_o$) and the $\lambda$-table (past steps), using four constants of the placed contact. Hence, by Theorem~main of the 11-22 c2 tex, \textbf{(Role) $\Leftarrow$ BS$(L)$ on the re-run-$b$ records only}. Case~b is the only open part of (Role).}
\author{Clio}\date{23 November 2026 (c2)}
\begin{document}\maketitle

\section*{Provenance}
\begin{itemize}
\item Code: \texttt{proofs/code-1123c2-Q597/} (new this session; no script shadows a stdlib module). Finite checks on the configuration table of Theorem~AW (\texttt{code-1120c2-Q597/words.py:configs}, through the \texttt{hsderive.py} preamble, as \texttt{caseA.py}): \texttt{grid.py} (Lemma~\ref{lem:grid} table), \texttt{dcomp.py} (Lemma~\ref{lem:dcomp}), \texttt{wakeA.py}, \texttt{wakeA2.py}, \texttt{wakeA3.py}, \texttt{boundbox.py}, \texttt{boundbox2.py} (Proposition~\ref{prop:wake}), and \texttt{lesionA.py} (the brief's lesion). Census corroboration: \texttt{nestdump.py} (the \texttt{rolesdump.py} loop with nest data added), \texttt{eaA.py}, \texttt{eaB.py} (Theorem~\ref{thm:ea} on data), and \texttt{staleA.py} (composition of $D$ on the 11-22 c2 pickles). DREAM probe: \texttt{dinvprobe.py}.
\item Upstream, all my own files, read first-hand this session: \texttt{2026-11-22c2-Q597-role.tex} (Theorem~main, Proposition~casea, Lemma~suffix, Lemma~term, ``(Stale), reduced''), \texttt{2026-11-21c2-Q597-sheet.tex} (Lemma~fin, Lemma~phi, Proposition~sh3 and its past-half method), \texttt{2026-11-22-Q597-first-contact.tex} (Lemma~dinv, Lemma~U), \texttt{2026-11-23-Q597-share.tex} (Lemma~shift, Lemma~mono), and the engine \texttt{code-1112c2-Q585/migrate.py} (\texttt{Board.journey}, \texttt{packed\_at}, \texttt{coef}, \texttt{is\_ideal}). No external source is used.
\item Brief question: \emph{does case~a need Lemma~U?} Proposition~casea does: Lemma~U forces the two second steps and drives Lemma~suffix. The three new results here (Theorem~\ref{thm:ea}, Lemma~\ref{lem:dcomp}, Proposition~\ref{prop:wake}) do not use it.
\end{itemize}

\section*{Summary}
\begin{enumerate}
\item \textbf{Theorem~\ref{thm:ea} ((Ea), proved).} The termination test at $R$ requires $f_0^T(\min R)=f_0^T(X_T)$ and $f_1^T(\min R)=f_1^T(X_T)$ (Lemma~\ref{lem:grid}, read from the engine). Here $f_0^B=a_0+b_0$, $f_1^B=a_1+b_1$, $f_0^C=b_0$ and $f_1^C=b_1$. Along $T$-journeys both functionals are monotone (non-decreasing for $B$, non-increasing for $C$). $f_0^T$ moves strictly at $(1,a0)$ and $f_1^T$ at $(1,a1)$ (finite table). In case~a, $A$ and $A'$ leave $R_p$ by $(1,a_T)$ and $(1,a_{T'})$. Each of these moves strictly the functional that the other one fixes. So neither journey can stop at its next rhombus. On the census, the test never fires there (258/258), and $R_p$ is off the grid in \emph{both} blocks in every record of either case (779/779). Control: the re-implemented test fires at the end of all $3\times1558$ journeys and never earlier.
\item \textbf{Lemma~\ref{lem:dcomp} (proved, one finite check).} $D=X_{p+2}\setminus Y_{p+2}$ consists of wake tiles and of non-rhombus tiles created by $A$ at $p$ or $p+1$. The 11-22 c2 kind (M$'$) is absent in case~a. The re-run's two hexagons meet, outside $A$'s two, only $J_o$-created tiles: 6 of 6 in all 8 placements. Lesion: with $J_o$'s \emph{removed} tiles in place of the created ones, the claim fails in 8/8. Census: 258/258 and 348/348 records have only these two kinds.
\item \textbf{Proposition~\ref{prop:wake} ((W$_a$), proved by hand from placement constants).} $A$ reads no wake tile at any step $k\ge p+2$. For $J_o$'s steps $j\ge c_0$ ($T=B$), resp.\ $j\ge c_0+2$ ($T=C$), use Lemma~phi: one functional, monotone in opposite senses on the two targets. For $j<c_0$, $\lambda$ is non-decreasing on both journeys, and removed tiles sit at $\lambda$-offset $\ge0$ and $\le\max Q_{T'}$. One coordinate of $\lambda(R_{p+2})-\lambda(R^J_{c_0})$ equals $1$ where $\max Q_{T'}$ is $0$. For $T=C$, $j\in\{c_0,c_0+1\}$, one explicit triangle remains, and its $\lambda_1$ lies below $\lambda_1(R_{p+2})$. Controls: shpast's common-rhombus window is non-empty (48 cells, the 21 tiles), and the lesioned references give 159 non-empty cells.
\item \textbf{Theorem~\ref{thm:casea}.} (Stale$_a$) holds. With Proposition~casea, BS$(L)$ and (Ca) hold in case~a. Theorem~main of 11-22 c2 now reads \textbf{(Role) $\Leftarrow$ BS$(L)$ restricted to the records whose $L$-re-run takes $b$}.
\item \textbf{Brief lesion} (drop the survivor clause). (Ea) is survivor-free by design: it uses only the two step-$p$ types. The (Stale$_a$) ingredients use the placement, and on the shifted offsets they fail (\S\ref{sec:lesion}).
\item \textbf{DREAM probe} (Lemma~dinv for $\Pi_C-\Pi_B$): negative. $\Pi_C-\Pi_B$ is a function of the rhombus shape alone. It is $0$, $(0,0,-1,0)$ or $(0,0,0,1)$ for the nest-$A$, $B$-class-1 and $C$-class-1 shapes, one value per shape on every rhombus of the census. It changes exactly at b-steps (891 of 1558 prefixes) and carries only the class, so it is no tool for (Resync).
\end{enumerate}

\section{Setting}\label{sec:set}
As in \S1 of the 11-22 c2 tex. $(x,y)$ is a diverging $B/C$ pair, and $L$ is its role at a class-1 a-step $(1,a_T)$ of its own target $T$ at its contact $p$ (Proposition~orient there). $A=J_L$ toward $T$ runs on $M_0$, with boards $X_k$ and rhombi $R_k$. Its re-run $A'$ toward $T$ runs on $M_x=M_0\cdot J_o$, with boards $Y_k$ and rhombi $R'_k$. $J_o$ (toward $T'\ne T$) runs on $M_0$ with rhombi $R^J_j$, and $W=M_0\setminus M_x$. Write $a_B=a1$, $a_C=a0$. \textbf{Case~a}: $A'[p]$ is $(1,a_{T'})$. By Theorem~local (first-contact tex) the contact is then one of four placed survivors per target. $A$ is at $(1,a_T)$, the re-run at $(1,a_{T'})$, and $|H_p\cap W|=3$, with $J_o$'s steps $c_0,c_0+1,c_0+2$ read. The two survivor rows per target and the two admissible types of step $c_0+2$ are those of \texttt{caseA.py}. $\alpha$ and $\alpha'$ are the rhombus displacements of $A[p]$ and $A'[p]$. By the definition of the types, $\alpha\in\Lambda_s$ and $\alpha'\in\Lambda_{1-s}$, with $s=1$ for $T=B$ and $s=0$ for $T=C$.

Proposition~casea (11-22 c2) proves: \emph{if} (Ea) and (Stale$_a$) hold, then $L$'s middle is $(1,a_T)(1,a_{T'})$, the re-run's is its translate-swap, and BS$(L)$ and (Ca) hold. Its proof uses Theorem~local, Lemma~U, Lemma~suffix and Lemma~term, all proved. This file proves the two hypotheses:
\begin{itemize}
\item[\textbf{(Ea)}] neither $A$ nor $A'$ stops at $R_{p+1}$, resp.\ $R'_{p+1}$;
\item[\textbf{(Stale$_a$)}] no hexagon $H_k$ of $A$ with $k\ge p+2$ meets $D:=X_{p+2}\setminus Y_{p+2}$.
\end{itemize}
In (Stale$_a$) the hexagons are those of Lemma~suffix (i) at $k=p+2$. That lemma is applied only once both journeys stand on $R_{p+2}=R'_{p+2}$, which Proposition~casea derives from (Ea).

\section{(Ea): the stop grid}\label{sec:ea}
\begin{lemma}[the stop grid]\label{lem:grid}
Let a journey toward $T\in\{B,C\}$ stop at the rhombus $R$. Then
\[f_0^T(\min R)=f_0^T(X_T),\qquad f_1^T(\min R)=f_1^T(X_T),\]
where $X_T$ is nest $T$'s frame origin, $f_0^T$ is the linear form on $\Lambda_0$ with kernel $\mathbb{R}u_T$, and $f_1^T$ is the linear form on $\Lambda_1$ with kernel $\mathbb{R}w_T$ (each extended by $0$ on the other lattice). With the engine's frames $u_B=(-1,1;0,0)$, $w_B=(0,0;1,-1)$, $u_C=(1,0;0,0)$, $w_C=(0,0;-1,0)$ (Lemma~fin of the sheet tex), these forms are $f_0^B=a_0+b_0$, $f_1^B=a_1+b_1$, $f_0^C=b_0$ and $f_1^C=b_1$.
\end{lemma}
\begin{proof}
By Lemma~fin (sheet tex) a journey toward $T$ stops only for nest $T$. Its rhombi have the nest-$A$ class (the source, skipped) or $T$'s class~1, and \texttt{packed\_at}$(T',R)$ is \texttt{None} at both (as in Lemma~term). \texttt{packed\_at}$(T,R)$ (read first-hand in \texttt{migrate.py}) returns a value only if some vertex $q$ of $R$ has $R=\{q,q+u_T,q+w_T,q+u_T+w_T\}$. Then $D=q-X_T$ is split as $D=D_u+D_w$ into its $\Lambda_0$ and $\Lambda_1$ blocks ($u_T\in\Lambda_0$, $w_T\in\Lambda_1$), and \texttt{coef} must return numbers $i,j$ with $D_u=iu_T$ and $D_w=jw_T$. Otherwise the vertex is skipped. So $f_0^T(D)=f_0^T(D_u)=i\,f_0^T(u_T)=0$, and likewise $f_1^T(D)=0$. Finally $\min R-q\in\{0,u_T,w_T,u_T+w_T\}$, which both forms annihilate. So $f_\epsilon^T(\min R)=f_\epsilon^T(q)=f_\epsilon^T(X_T)$.
\end{proof}
\noindent \emph{Engine control} (\texttt{eaA.py}): the test, re-implemented from \texttt{migrate.py}, fires at the final rhombus of every $A$, re-run and $J_o$ of all 1558 ordered size-5 pairs, and at no earlier rhombus.

\begin{lemma}[the grid functionals along a journey]\label{lem:gridmono}
Put $F_\epsilon^T(R)=f_\epsilon^T(\min R)$. Along every step of a terminating journey toward $T$, the change $\Delta F_\epsilon^T$ is:
\begin{center}\small\begin{tabular}{lcccccc}\toprule
& $(0,a0)$ & $(0,a1)$ & $(0,b)$ & $(1,a0)$ & $(1,a1)$ & $(1,b)$\\\midrule
$T=B$: $\Delta F_0^B$, $\Delta F_1^B$ & $0,0$ & $0,1$ & $1,1$ & $\mathbf{1},0$ & $0,\mathbf{1}$ & $0,0$\\
$T=C$: $\Delta F_0^C$, $\Delta F_1^C$ & $-1,0$ & $0,0$ & $0,0$ & $\mathbf{-1},0$ & $0,\mathbf{-1}$ & $-1,-1$\\\bottomrule
\end{tabular}\end{center}
So $F_0^B$, $F_1^B$ are non-decreasing along $B$-journeys, and $F_0^C$, $F_1^C$ are non-increasing along $C$-journeys. At the class-1 a-steps, $(1,a0)$ moves $F_0^T$ strictly and fixes $F_1^T$, and $(1,a1)$ moves $F_1^T$ strictly and fixes $F_0^T$.
\end{lemma}
\begin{proof}[Proof (finite check)]
By Lemma~word every step is a translate of one of the six good configurations of its target. $\Delta F$ is translation invariant, because $\min(R+v)=\min R+v$ and $f$ is linear. \texttt{grid.py} reads the twelve configurations (\texttt{grid.out}). The row $F_0^B$ is Lemma~phi of the sheet tex ($F_0^B=g$). The two class-1 a-step entries also follow by hand. $\alpha$ of $(1,a1)$ lies in $\Lambda_1$, so $F_0$ is fixed, and $\alpha$ of $(1,a0)$ lies in $\Lambda_0$, so $F_1$ is fixed.
\emph{Lesion:} $C$'s functionals along $B$-journeys take $\Delta\in\{-1,0,1\}$, so they are not monotone (\texttt{grid.out}). The pairing of each functional with its own target matters.
\end{proof}

\begin{theorem}[(Ea)]\label{thm:ea}
In case~a, $A$ does not stop at $R_{p+1}$ and $A'$ does not stop at $R'_{p+1}$.
\end{theorem}
\begin{proof}
Let $T=B$. $A$ and $A'$ are terminating journeys toward $B$ with the same nest~$B$: $J_o$ moves toward $C$ and changes only the positions of nests $A$ and $C$, and the frame origin $X_B$ is fixed. Let $R^{\rm end}$ and $R'^{\rm end}$ be their final rhombi. By Lemma~\ref{lem:grid}, $F_\epsilon^B(R^{\rm end})=F_\epsilon^B(R'^{\rm end})=f_\epsilon^B(X_B)=:c_\epsilon$.
$A'$ takes $(1,a0)$ at $R_p$ (case~a), so $F_0^B(R'_{p+1})=F_0^B(R_p)+1$. By monotonicity (Lemma~\ref{lem:gridmono}) $c_0=F_0^B(R'^{\rm end})\ge F_0^B(R_p)+1$. $A$ takes $(1,a1)$ at $R_p$, which fixes $F_0^B$, so $F_0^B(R_{p+1})=F_0^B(R_p)<c_0$, and by Lemma~\ref{lem:grid} $A$ does not stop at $R_{p+1}$.
Symmetrically, $A$'s step $(1,a1)$ gives $c_1=F_1^B(R^{\rm end})\ge F_1^B(R_p)+1$, while $A'$'s step $(1,a0)$ fixes $F_1^B$. So $F_1^B(R'_{p+1})=F_1^B(R_p)<c_1$, and $A'$ does not stop at $R'_{p+1}$.
For $T=C$, all inequalities reverse. $A$ takes $(1,a0)$ and $A'$ takes $(1,a1)$. $A'$ gives $c_1\le F_1^C(R_p)-1$, while $F_1^C(R_{p+1})=F_1^C(R_p)$. $A$ gives $c_0\le F_0^C(R_p)-1$, while $F_0^C(R'_{p+1})=F_0^C(R_p)$.
The argument uses only the two step-$p$ types, which hold by the definition of case~a, and not the survivor table. Neither journey's later behaviour enters, so there is no circularity.
\end{proof}
\noindent\emph{Data} (\texttt{eaA.py}, \texttt{eaB.py}, size~5 full census, nest data from \texttt{nestdump.py}): the test fires at neither $R_{p+1}$ nor $R'_{p+1}$ in 258/258 case-a records. $R_p$ is off the grid in both blocks, $f_0$ and $f_1$, in all 258 case-a and all 521 case-b records. $R_{p+1}$ ($T=B$) keeps $R_p$'s off-grid $\Lambda_0$ block, and $R'_{p+1}$ keeps its $\Lambda_1$ block, in every record, as the proof predicts. The size-6 subsample is reported in \S\ref{sec:six}.

\section{(Stale$_a$)}\label{sec:stale}
\begin{lemma}[composition of $D$]\label{lem:dcomp}
In case~a, assume (Ea). Then every tile of $D=X_{p+2}\setminus Y_{p+2}$ is either a wake tile ($\in W$) or a non-rhombus tile created by $A$ at step $p$ or $p+1$.
\end{lemma}
\begin{proof}
By Proposition~casea's first two paragraphs (which need only (Ea), Lemma~U and the placement), $A[p+1]$ is $a_{T'}$ and $A'[p+1]$ is $a_T$, with $R_{p+2}=R'_{p+2}$. Write $H_p,H_{p+1}$ for $A$'s hexagons, $H'_p,H'_{p+1}$ for $A'$'s, and $C_p,C_{p+1},C'_p,C'_{p+1}$ for the created tiles. Then
\[X_{p+2}\subseteq\bigl(X_p\setminus(H_p\cup H_{p+1})\bigr)\cup C_p\cup C_{p+1},\qquad Y_{p+2}\supseteq Y_p\setminus(H'_p\cup H'_{p+1}).\]
Let $t\in D$.
\emph{(a) $t\in C_p\cup C_{p+1}$.} The rhombi there are $R_{p+1}$, which $A$ removes at $p+1$, and $R_{p+2}=R'_{p+2}\in Y_{p+2}$. So $t$ is a non-rhombus tile created by $A$ at $p$ or $p+1$.
\emph{(b) $t\in X_p\setminus(H_p\cup H_{p+1})$.} If $t\notin Y_p$, then $t\in X_p\setminus Y_p=M_0\setminus M_x=W$ by Lemma~dinv. If $t\in Y_p$, then $t\notin Y_{p+2}$ forces $t\in H'_p\cup H'_{p+1}$. It is not a rhombus: $H'_p$'s rhombus is $R_p\in H_p$, and $H'_{p+1}$'s rhombus $R'_{p+1}$ was created at $p$, so it is not in $Y_p$. It is not in $C'_p$, because a tile already present is not created again. By the finite check below, every non-rhombus tile of $(H'_p\cup H'_{p+1})\setminus(H_p\cup H_{p+1})$ lies in $J_o$'s created set $\mathrm{Jcre}$ (its steps $c_0,c_0+1,c_0+2$) or in $C'_p$. So $t\in\mathrm{Jcre}$. But $\mathrm{Jcre}\subseteq M_x\setminus M_0=Y_p\setminus X_p$: $J_o$ never removes a non-rhombus tile it made (no-revisit), never re-creates a removed one (Lemma~nr), and a present tile is not created; the equality is Lemma~dinv. This contradicts $t\in X_p$. So case~(b) gives only wake tiles.
\end{proof}
\noindent\emph{Finite check} (\texttt{dcomp.py}, the placement of \texttt{caseA.py}): in all 8 placements (4 survivors $\times$ 2 admissible types of $J_o$'s step $c_0+2$), the re-run's two hexagons have exactly 6 non-rhombus tiles outside $H_p\cup H_{p+1}$, and all 6 are in $\mathrm{Jcre}$. $H'_{p+1}\cap C'_p=\emptyset$ (no-revisit, as it must be), and $R_{p+2}=R'_{p+2}$. \emph{Lesion:} with $J_o$'s removed tiles instead of $\mathrm{Jcre}$ the inclusion fails in 8/8.
\emph{Census} (\texttt{staleA.py}, 11-22 c2 pickles): $D$ has exactly the two kinds in 258/258 (size 5) and 348/348 (subsample) records, with 3 tiles created at $p$ and 3 at $p+1$ in every record. No tile of either kind is in a later hexagon of $A$.

\begin{proposition}[(W$_a$): no wake read after $p+1$]\label{prop:wake}
In case~a, assume (Ea). Then no hexagon $H_k$ of $A$ with $k\ge p+2$ contains a tile of $W$.
\end{proposition}
\begin{proof}
$H_k$'s rhombus $R_k$ ($k\ge1$) was created by $A$, so it is not in $M_0\supseteq W$. Let $t\in H_k\cap W$ be a non-rhombus tile. By no-revisit and Lemma~nr (as in Lemma~hs of 11-22 c2), $J_o$ removes $t$ at exactly one step $j$. Two facts are used. \textbf{(LAM)}: $\lambda=(\min a_0,\min a_1)$ is componentwise non-decreasing along every journey (G1c tex table, Lemma~word). \textbf{($Q_\lambda$)}: a non-rhombus tile $t$ removed at a step from $R$ toward $T$ has $\lambda(t)-\lambda(R)\in\{(0,0),(1,0)\}$ for $T=B$ and $\in\{(0,0),(0,1)\}$ for $T=C$. These are the $\lambda$-parts of $Q_T$ in Lemma~window, re-read from the twelve configurations by \texttt{wakeA.py}. Also needed are four constants of the placed contact, read off by \texttt{wakeA2.py} and \texttt{boundbox2.py}. The $g$-columns are the same for both survivors and both types of step $c_0+2$. The $\lambda$-column takes the value shown first for one survivor and the second for the other; only the coordinate used below matters, and it is $1$ in both:
\begin{center}\small\begin{tabular}{lcccc}\toprule
& $g(R_p)-g(R^J_{c_0})$ & $g(R_{p+2})-g(R_p)$ & $g(R^J_{c_0+2})-g(R_p)$ & $\lambda(R_{p+2})-\lambda(R^J_{c_0})$\\\midrule
$T=B$ & $0$ & $+1$ & (not used) & $(1,0)$ or $(1,1)$\\
$T=C$ & $+1$ & $-1$ & $0$ & $(0,1)$ or $(1,1)$\\\bottomrule
\end{tabular}\end{center}
Here $g(R)=\varphi(\min R)$ and $\varphi=a_0+b_0$ (Lemma~phi). The second column also follows from Lemma~phi's table: $(1,a1)$ then $(1,a0)$ toward $B$ gives $0+1$, and $(1,a0)$ then $(1,a1)$ toward $C$ gives $-1+0$.
\emph{Past, $j<c_0$.} Let $T=B$. $A$'s tile: $\lambda_0(t)\ge\lambda_0(R_k)\ge\lambda_0(R_{p+2})=\lambda_0(R^J_{c_0})+1$, by ($Q_\lambda$) and (LAM). $J_o$'s tile ($J_o$ toward $C$): $\lambda_0(t)\le\lambda_0(R^J_j)+0\le\lambda_0(R^J_{c_0})$. Contradiction. For $T=C$, the same with $\lambda_1$: $A$ gives $\lambda_1(t)\ge\lambda_1(R^J_{c_0})+1$, and $J_o$ (toward $B$, whose $\lambda_1$-offset is $0$) gives $\lambda_1(t)\le\lambda_1(R^J_{c_0})$.
\emph{Future and contact steps, $T=B$, $j\ge c_0$.} By Lemma~phi (ii), $A$'s tile has $\varphi(t)=g(R_k)\ge g(R_{p+2})=g(R_p)+1$. $J_o$'s tile (toward $C$) has $\varphi(t)\le g(R^J_j)\le g(R^J_{c_0})=g(R_p)$. Contradiction.
\emph{Future, $T=C$, $j\ge c_0+2$.} $A$'s tile (toward $C$) has $\varphi(t)\le g(R_k)\le g(R_{p+2})=g(R_p)-1$. $J_o$'s tile (toward $B$) has $\varphi(t)=g(R^J_j)\ge g(R^J_{c_0+2})=g(R_p)$. Contradiction.
\emph{Contact steps, $T=C$, $j\in\{c_0,c_0+1\}$.} These two hexagons are placed explicitly. A tile of them that $A$ removes at $p$ or $p+1$ is not removed again. Exactly one non-rhombus tile is not removed there, a triangle, in each placement (\texttt{wakeA3.py}). It has $\lambda(t)-\lambda(R_{p+2})=(0,-1)$, resp.\ $(-1,-1)$. Any tile $A$ removes at $k\ge p+2$ has $\lambda(t)\ge\lambda(R_k)\ge\lambda(R_{p+2})$ by ($Q_\lambda$) and (LAM). So $A$ never removes it.
\end{proof}
\noindent\emph{Cross-checks.}
\begin{itemize}
\item The past half was first run as shpast's window-and-words check, with reference rhombi $R_{p+2}$ for $A$ and $R^J_{c_0}$ for $J_o$ (\texttt{wakeA.py}, \texttt{wakeA2.py}). It found 0 feasible tiles at boxes 4, 6 and 8 (\texttt{wakeA4.out}, \texttt{wakeA2.out}, \texttt{wakeA2-box8.out}), while the lesioned references give 174 feasible tiles at both box 6 and box 8, none on the boundary. Exact interval propagation over all 7 shapes and all $(q,q')$ (\texttt{boundbox.py}, \texttt{boundbox2.py}) shows that the $\lambda$-part alone empties all 175 cells. That is the hand proof above, and no box is involved.
\item \emph{Controls.} With shpast's common class-0 rhombus $R^\ast$, the same instrument finds 48 non-empty cells, containing shpast's 21 tiles (\texttt{wctl.py} reproduces the 21). \emph{Lesion:} with references $R_p$ (so $A$'s steps $p,p+1$ count) and $R^J_{c_0+3}$ (so the read steps count as past), there are 159 non-empty cells. The one-sided window of the triangle step is live: the tiles $A$ removes at $p+1$ pass it from $R_{p+1}$ in 8/8.
\item \emph{Census:} $A$ reads no wake tile after $p+1$ in 258/258 and 348/348 records (\texttt{caseAdata*.out}, 11-22 c2), and does read one at $p+1$ in all of them.
\end{itemize}

\begin{theorem}[case a closes]\label{thm:casea}
Let $(x,y)$ be a diverging $B/C$ pair satisfying (FC$\le$) on both roles, and let $L$'s re-run take $(1,a_{T'})$. Then (Ea) and (Stale$_a$) hold. Hence (Proposition~casea, 11-22 c2) $L$'s middle is $a\cdot c$ with $c=A'[p]+\alpha$ a $(1,a_{T'})$-step, the re-run's middle is $(c-\alpha)(a+\alpha')$, and BS$(L)$ and (Ca) hold. By Theorem~main of 11-22 c2, $I=\emptyset$ for every such pair.
\end{theorem}
\begin{proof}
(Ea) is Theorem~\ref{thm:ea}. For (Stale$_a$): by Lemma~\ref{lem:dcomp}, a tile of $D$ is a wake tile or a non-rhombus tile made by $A$ at $p$ or $p+1$. No $H_k$ with $k\ge p+2$ contains a wake tile (Proposition~\ref{prop:wake}). None contains a non-rhombus tile made by $A$, by no-revisit (no-repeat tex, Theorem~main). (Share), needed by Proposition~casea through Proposition~orient, is proved (share tex, 11-23).
\end{proof}

\begin{corollary}[(Role) needs only case b]\label{cor:role}
In Theorem~main of the 11-22 c2 tex, hypothesis (Ca) and the case-a half of BS$(L)$ are now theorems. So \textup{(Role)} $\Leftarrow$ \textup{BS}$(L)$ for the pairs whose $L$-re-run takes $(1,b)$, which are the geometric records of Remark~reindex.
\end{corollary}

\section{The brief's lesion: drop the survivor clause}\label{sec:lesion}
\begin{itemize}
\item \textbf{(Ea)} does not use the survivor table, by design. It uses the step-$p$ types $(1,a_T)$ and $(1,a_{T'})$ only, which define case~a. So the lesion cannot touch it, and what it does use is named: the stop grid (Lemma~\ref{lem:grid}) and the monotone table (Lemma~\ref{lem:gridmono}). The table has its own lesion, the wrong target's functionals, which are not monotone.
\item \textbf{(Stale$_a$)} uses the placement through Lemma~\ref{lem:dcomp} (which re-run tiles are $J_o$-created) and through the four constants of Proposition~\ref{prop:wake}. \texttt{lesionA.py} evaluates all three ingredients (DC, FR, P) on the $|S|=2$ placements and on the 12 unit shifts of each survivor offset. Result (\texttt{lesionA.out}):
\begin{center}\small\begin{tabular}{lrrrrr}\toprule
placements ($T=B$ / $T=C$) & number & DC & FR & P empty & all three\\\midrule
true survivors, $|S|=3$ & 4 / 4 & 4 / 4 & 4 / 4 & 4 / 4 & \textbf{4 / 4}\\
true offset, $|S|=2$ (non-survivor $c_0+2$ type) & 2 / 2 & 0 / 0 & 2 / 2 & 2 / 2 & 0 / 0\\
shifted offsets (placeable) & 72 / 72 & 0 / 0 & 32 / 36 & 60 / 60 & 0 / 0\\\bottomrule
\end{tabular}\end{center}
So the three ingredients hold together exactly on the 8 survivor placements. The survivor clause enters through Lemma~\ref{lem:dcomp}: DC fails on every lesioned placement. FR and P alone are less selective. That is expected, since they compare only potentials of two rhombi, so they survive many shifts, which is why the $D$-composition is needed.
\end{itemize}

\section{DREAM probe: is $\Pi_C-\Pi_B$ invariant along a common prefix?}
No, and it could not be a tool. By Lemma~shift of the share tex, $\Pi_B$ and $\Pi_C$ translate by the same vector, so $\Pi_C(R)-\Pi_B(R)$ depends only on the shape of $R$. \texttt{dinvprobe.py} confirms this on every rhombus of every journey in \texttt{roles5.pkl}: three shapes, one value each, $(0,0,0,0)$ (nest-$A$ shape), $(0,0,-1,0)$ ($B$ class~1) and $(0,0,0,1)$ ($C$ class~1). Along a common prefix it changes exactly at b-steps (891 of 1558 prefixes change, and every change follows a b-step). It records the class and nothing else, so it carries no information for (Resync). I did not pursue it further, as the brief asked.

\section{Retreat probe (case b): (Comp)~(i) reduces to a statement about $J_o$ at the split}\label{sec:compi}
Case~a closed early, so the brief's retreat target was probed as well, time-boxed and labelled as a probe. In case~b, (Comp)~(i) says that $H_p+\gamma$ contains no tile of $M_0\cap M_x$. Here $\gamma:=R_{p+m+1}-R_{p+1}$ is defined by $A$'s own middle, and $R_{p+m+1}=R^\sharp$ is $A$'s image of the split rhombus $R^\ast=R_{p+m}=R^o_e$. So $H_p+\gamma$ is the $(1,a_T)$ configuration at $R^\sharp-\alpha$, and this needs no BS.
\begin{itemize}
\item \emph{Data} (\texttt{lastsplit.py}): the re-run's last middle hexagon has a square and a triangle created by $J_o$ at $e$, and a triangle created by $J_o$ at $e+1$. This holds in 521/521 and 506/506 records. $J_o$'s step $e$ is $(0,b)$ in all of them.
\item \emph{Finite check} (\texttt{compi.py}, at the shared split hexagon): the square and one triangle of the $(1,a_T)$ configuration at $R^\sharp-\alpha$ are created by $J_o$'s $(0,b)$ at $R^\ast$. The third triangle is created by $J_o$'s step $e+1$ exactly when that step is $(1,a_T)$ or $(1,b)$, and not when it is $(1,a_{T'})$. \emph{Lesions:} with $\alpha'$ or $0$ in place of $\alpha$, no tile is $J_o$-created.
\item Hence \textbf{(Comp)~(i) $\Leftarrow$ (E1)}: \emph{$J_o$ takes a step at $R^o_{e+1}$, and it is not $(1,a_{T'})$.} $J_o$-created tiles are not in $M_0$ (Lemma~nr). (E1) is computed: $J_o$'s step $e+1$ is $(1,a_T)$ or $(1,b)$ in 521/521 and 506/506 records (\texttt{e1.out}).
\item (E1) is \emph{not} a word property. Over all journeys of the census, the factor ``$(0,b)$ then the target's own $(1,a)$'' occurs 1346 times at size 5, and a journey ends right after a $(0,b)$ 443 times (\texttt{factor.out}). It is not forced by no-revisit either: no class-1 configuration at $R^o_{e+1}$ meets $J_o$'s own tiles from step $e$ (\texttt{e1.py}). So (E1) is a genuine two-journey statement at the split. A likely route is certification by $A$'s next hexagon at $R^\sharp$, whose non-rhombus tiles lie in $M_0\cap M_x$ (Sh3 and no-revisit) and so are present on $J_o$'s board. That route is unavailable when the split is $A$'s last step (180/521). Not pursued further.
\end{itemize}

\section{Census corroboration and the size-6 subsample}\label{sec:six}
\begin{center}\small\begin{tabular}{lrr}\toprule
& size 5 (full, 5194 states) & size 6 (\texttt{KEEP=3}, 3951 states)\\\midrule
pairs dumped by \texttt{nestdump.py} (diverging $B/C$) & 779 (3\,min\,59\,s) & 854 (6\,min\,1\,s)\\
CONTROL: test fires at the end of $A$, re-run, $J_o$ / earlier & $3\times1558$ / 0 & $3\times1708$ / 0\\
case-a records ($T=B$, $T=C$) & 211, 47 & 295, 53\\
(Ea): test fires at $R_{p+1}$ or $R'_{p+1}$ & 0 & 0\\
$R_p$ off-grid in both blocks, case a + case b & 258 + 521 & 348 + 506\\
$R_{p+1}$ keeps $R_p$'s $s$-complement block, $R'_{p+1}$ keeps the other & all & all\\
$D$ = wake $\cup$ ($A$-made at $p,p+1$) (\texttt{staleA.py}) & 258/258 & 348/348\\
$A$ reads a wake tile after $p+1$ (\texttt{caseAdata.py}) & 0 & 0\\\bottomrule
\end{tabular}\end{center}
A consistency check on the grid picture: at size~5, $R_{p+2}$ is on the grid in both blocks in exactly 28 records, all with $T=C$, and these are exactly the 28 records in which $A$ ends at $p+2$ (\texttt{staleA5.out}). On the subsample both counts are 19. For $T=B$, $R_{p+2}$ is never on the $\Lambda_1$ line, and $A$ never ends at $p+2$.

\section{Status and what remains}
\begin{itemize}
\item \textbf{Proved this session:} Lemma~\ref{lem:grid} (by hand, from the engine); Lemma~\ref{lem:gridmono} (finite check on the twelve configurations); Theorem~\ref{thm:ea}; Lemma~\ref{lem:dcomp} (one finite check on the 8 placements); Proposition~\ref{prop:wake} (by hand, from Lemma~phi, (LAM), ($Q_\lambda$) and four placement constants); Theorem~\ref{thm:casea}; Corollary~\ref{cor:role}.
\item \textbf{Inputs, all previously proved:} Theorem~local and the \texttt{caseA.py} placement; Lemma~U; Lemma~dinv; Lemma~suffix; Lemma~term; Lemma~fin; Lemma~phi; Lemma~word and the G1c $\Delta\lambda$ table; no-revisit and Lemma~nr; (Share).
\item \textbf{Open:} (Role) now waits only on BS$(L)$ for the case-b (geometric) records, i.e.\ on the non-circular route of Remark~order (11-22 c2): (A3), (Tn), WS$_{\rm last}$, (Comp). The retreat target (Comp)~(i) is reduced to (E1), a statement about $J_o$'s step after the split (\S\ref{sec:compi}, computed, not proved).
\item The line-by-line $I$-freeness re-read of the post-G1a stages (Remark~reindex, 11-22 c2) is still owed.
\end{itemize}
\end{document}
